\subsection{Hypergradient Descent}

% ---------------------------------------------------------------
% 1. Explicación gráfica
% ---------------------------------------------------------------
\begin{frame}{Hypergradient Descent: Explicación Gráfica}
    \begin{block}{Concepto principal}
        \scriptsize
        La tasa de aprendizaje $\alpha$ es el hiperparámetro crítico de todo método de
        primer orden. La idea: tratarla como \emph{una variable más} y aplicarle
        \textbf{descenso de gradiente a ella misma}.
        \vspace{-0.3em}
        \[
            \frac{\partial f(\mathbf{x}^{(k)})}{\partial \alpha}
            = \nabla f(\mathbf{x}^{(k)})^{\top}
              \frac{\partial \mathbf{x}^{(k)}}{\partial \alpha}
            = -\,\nabla f(\mathbf{x}^{(k)})^{\top} \nabla f(\mathbf{x}^{(k-1)})
        \]
        \vspace{-0.5em}

        El \textbf{hipergradiente} es el producto punto de los \textbf{dos últimos
        gradientes}: información que ya se tiene, gratis.
    \end{block}

    \vspace{-0.2em}
    \begin{center}
    \begin{tikzpicture}[x=0.42cm, y=0.33cm, font=\tiny]
        % ---- panel izquierdo: gradientes alineados ----
        \draw[CelesteBrillante, thick, domain=-2.4:2.4, samples=40, smooth]
              plot (\x, {0.45*\x*\x});
        \fill[RojoAlerta] (2.2,2.178) circle (1.4pt);
        \fill[RojoAlerta] (1.4,0.882) circle (1.4pt);
        \draw[NaranjaBase, ->, very thick] (2.2,3.4) -- (1.0,3.4);
        \draw[NaranjaBase, ->, very thick] (1.4,4.4) -- (0.2,4.4);
        \node[above, text=VerdeNeon] at (0.9,4.8) {$g^{(k)}\!\cdot g^{(k-1)} > 0$};
        \node[below, text=GrisTexto, align=center] at (0,-0.4)
             {mismo sentido: voy bien\\ \textcolor{VerdeNeon}{$\alpha$ \textbf{aumenta}}};
        % ---- panel derecho: gradientes opuestos ----
        \begin{scope}[xshift=5.4cm]
        \draw[CelesteBrillante, thick, domain=-2.4:2.4, samples=40, smooth]
              plot (\x, {0.45*\x*\x});
        \fill[RojoAlerta] (2.2,2.178) circle (1.4pt);
        \fill[RojoAlerta] (-1.8,1.458) circle (1.4pt);
        \draw[NaranjaBase, ->, very thick] (2.2,3.4) -- (1.0,3.4);
        \draw[NaranjaBase, ->, very thick] (-1.8,4.4) -- (-0.6,4.4);
        \node[above, text=RojoAlerta] at (0.2,4.8) {$g^{(k)}\!\cdot g^{(k-1)} < 0$};
        \node[below, text=GrisTexto, align=center] at (0,-0.4)
             {me pasé, oscilo\\ \textcolor{RojoAlerta}{$\alpha$ \textbf{disminuye}}};
        \end{scope}
    \end{tikzpicture}
    \end{center}
\end{frame}

% ---------------------------------------------------------------
% 2. Pseudocódigo
% ---------------------------------------------------------------
\begin{frame}{Hypergradient Descent: Pseudocódigo}
    \scriptsize
    \begin{algorithmic}[1]
        \Require $f$, $\nabla f$, punto $\mathbf{x}$, tasa inicial $\alpha_0$,
                 \textbf{hipertasa} $\mu$
        \State $\alpha \gets \alpha_0$; \quad
               $\mathbf{g}_{\text{prev}} \gets \mathbf{0}$
               \Comment{en el paso 0 no hay gradiente anterior}
        \For{$k = 0$ \textbf{to} $k_{\max}$}
            \State $\mathbf{g} \gets \nabla f(\mathbf{x})$
            \State $\alpha \gets \alpha + \mu\,
                   \big(\mathbf{g}^{\top}\mathbf{g}_{\text{prev}}\big)$
                   \Comment{descenso sobre $\alpha$: $\alpha - \mu\,\partial f/\partial\alpha$}
            \State $\mathbf{g}_{\text{prev}} \gets \mathbf{g}$
            \State $\mathbf{x} \gets \mathbf{x} - \alpha\,\mathbf{g}$
                   \Comment{paso normal de descenso}
        \EndFor
        \State \Return $\mathbf{x}$
    \end{algorithmic}

    \vspace{0.5em}
    \begin{block}{Observaciones}
        \tiny
        \begin{itemize}
            \item El signo $+$ en la línea 4 \emph{es} un descenso: como
                  $\partial f/\partial\alpha = -\mathbf{g}^{\top}\mathbf{g}_{\text{prev}}$,
                  restarlo equivale a sumar el producto punto.
            \item Se cambia el problema de ajustar $\alpha$ por el de ajustar $\mu$,
                  pero el método es \textbf{mucho menos sensible} a $\mu$ que a $\alpha$.
            \item La misma idea se aplica sobre cualquier método de primer orden
                  (momentum, Nesterov, Adam): basta derivar respecto a su $\alpha$.
        \end{itemize}
    \end{block}
\end{frame}

% ---------------------------------------------------------------
% 3. Ejemplo paso a paso
% ---------------------------------------------------------------
\begin{frame}{Hypergradient Descent: Ejemplo Paso a Paso}
    \begin{exampleblock}{Planteamiento}
        \scriptsize
        $f(x)=x^{2}$, \ $f'(x)=2x$, \ $x^{(0)}=1$, \ $\mu=0.01$.
        Se prueban \textbf{dos tasas iniciales deliberadamente malas} para ver cómo el
        método las corrige solo.
    \end{exampleblock}

    \vspace{-0.2em}
    \begin{columns}[T]
        \begin{column}{0.49\textwidth}
            \begin{center}
            \tiny
            \textcolor{VerdeNeon}{\textbf{Caso A: $\alpha_0=0.10$ (muy chico)}}\\[0.2em]
            \begin{tabular}{c r r r}
                \toprule
                $k$ & $x^{(k)}$ & $g\cdot g_{\text{prev}}$ & $\alpha$ \\
                \midrule
                0 & $+1.0000$ & $0.000$ & $0.1000$ \\
                1 & $+0.8000$ & $3.200$ & $0.1320$ \\
                2 & $+0.5888$ & $1.884$ & $0.1508$ \\
                3 & $+0.4112$ & $0.968$ & $0.1605$ \\
                4 & $+0.2792$ & $0.459$ & $0.1651$ \\
                5 & $+0.1870$ & $0.209$ & $0.1672$ \\
                \bottomrule
            \end{tabular}
            \end{center}
            \vspace{0.2em}
            \tiny
            Los gradientes mantienen el signo $\Rightarrow$ producto $>0$
            $\Rightarrow$ $\alpha$ sube de $0.10$ a $0.167$: el método
            \textbf{acelera solo}.
        \end{column}
        \begin{column}{0.49\textwidth}
            \begin{center}
            \tiny
            \textcolor{RojoAlerta}{\textbf{Caso B: $\alpha_0=0.90$ (muy grande)}}\\[0.2em]
            \begin{tabular}{c r r r}
                \toprule
                $k$ & $x^{(k)}$ & $g\cdot g_{\text{prev}}$ & $\alpha$ \\
                \midrule
                0 & $+1.0000$ & $0.000$  & $0.9000$ \\
                1 & $-0.8000$ & $-3.200$ & $0.8680$ \\
                2 & $+0.5888$ & $-1.884$ & $0.8492$ \\
                3 & $-0.4112$ & $-0.968$ & $0.8395$ \\
                4 & $+0.2792$ & $-0.459$ & $0.8349$ \\
                5 & $-0.1870$ & $-0.209$ & $0.8328$ \\
                \bottomrule
            \end{tabular}
            \end{center}
            \vspace{0.2em}
            \tiny
            $x$ \textbf{alterna de signo} (sobrepasa el mínimo en cada paso)
            $\Rightarrow$ producto $<0$ $\Rightarrow$ $\alpha$ baja: el método
            \textbf{frena solo}.
        \end{column}
    \end{columns}

    \vspace{0.3em}
    \scriptsize
    En ambos casos $\alpha$ se mueve hacia la zona estable sin que nadie la ajuste
    a mano. Para $f=x^{2}$ el paso ideal es $\alpha=0.5$.
\end{frame}

% ---------------------------------------------------------------
% 4. Código
% ---------------------------------------------------------------
\begin{frame}[fragile]{Hypergradient Descent: Código}
\begin{lstlisting}[language=Python]
import numpy as np

def hypergradient_descent(f, grad, x, alpha=0.1, mu=1e-2,
                          n=100, eps=1e-8):
    x = np.atleast_1d(np.asarray(x, float))
    g_prev = np.zeros_like(x)        # no hay gradiente anterior
    for _ in range(n):
        g = grad(x)
        alpha = alpha + mu * (g @ g_prev)   # paso sobre alpha
        if alpha < 0:                # salvaguarda practica
            alpha = 1e-6
        g_prev = g
        x = x - alpha * g            # paso sobre x
        if np.linalg.norm(g) < eps:
            break
    return x, alpha

f    = lambda v: v @ v               # f(x) = x^2
grad = lambda v: 2*v
# hypergradient_descent(f, grad, [1.0], alpha=0.1)
#   -> x ~ 0,  alpha subio de 0.10 a ~0.17
\end{lstlisting}
\end{frame}
